% =====================================================================
%  Chapitre 5 — Tests et validation
% =====================================================================
\chapter{Tests et validation}
\label{chap:tests}

\section{Stratégie de validation}

La validation s'appuie sur trois niveaux, choisis en fonction de ce qu'il s'agit de
prouver.

\textbf{Les services métier se testent sans serveur.} Les formules de readiness, de
qualification, de priorisation et de valorisation sont des fonctions sans effet de bord.
Elles se vérifient par comparaison de valeurs exactes, sans base de données ni requête
HTTP.

\textbf{Les règles d'autorisation se testent par l'API.} Un contrôle d'accès ne se
prouve pas en lisant le code de l'interface. Les tests d'autorisation construisent un
jeton pour un rôle donné, appellent l'endpoint directement et vérifient le code de
retour. C'est la seule manière de garantir qu'un utilisateur qui contourne l'interface
se heurte bien au serveur.

\textbf{Les parcours d'écran se vérifient dans un navigateur réel.} Les enchaînements
d'interface — connexion selon le rôle, ouverture des panneaux latéraux, recalcul de la
priorisation, export CSV — ont été exécutés dans un navigateur automatisé, avec capture
à chaque étape.

\section{Composition de la suite automatisée}

La suite compte \textbf{112~tests} répartis en sept fichiers.

\begin{table}[H]
\centering
\caption{Répartition de la suite de tests}
\label{tab:repartition-tests}
\small
\begin{tabular}{@{}L{5.2cm} C{1.6cm} L{7.2cm}@{}}
\toprule
\textbf{Fichier} & \textbf{Tests} & \textbf{Objet} \\
\midrule
\code{test\_permissions.py}        & 40 & Comportement de chacun des 11 rôles \\
\code{test\_login.py}              & 19 & Authentification et aide de développement \\
\code{test\_opportunities.py}      & 15 & Cycle de vie, qualification, priorisation \\
\code{test\_tenant\_isolation.py}  & 13 & Étanchéité entre organisations \\
\code{test\_diagnostic.py}         & 11 & Sémantique du readiness \\
\code{test\_scope.py}              & 10 & Bornage par département \\
\code{test\_endpoint\_protection.py}&  4 & Garde-fous structurels sur l'API \\
\midrule
\textbf{Total}                     & \textbf{112} & \\
\bottomrule
\end{tabular}
\end{table}

\section{Résultats}

\begin{table}[H]
\centering
\caption{Résultat de l'exécution de la suite}
\label{tab:resultats-tests}
\small
\begin{tabular}{@{}L{5.6cm} L{5.4cm} C{2.8cm}@{}}
\toprule
\textbf{Vérification} & \textbf{Objectif} & \textbf{Résultat} \\
\midrule
Suite backend complète & 112 tests & \fait \ 112 réussis \\
Durée d'exécution & Boucle de retour courte & 12,15 s \\
Typage du frontend & \code{tsc --noEmit} sans erreur & \fait \ aucune erreur \\
\bottomrule
\end{tabular}
\end{table}

\section{Détail des vérifications}

\subsection{Sémantique du readiness}

\begin{longtable}{@{}L{7.4cm} L{5.2cm} C{1.5cm}@{}}
\caption{Tests du calcul de readiness}
\label{tab:tests-readiness}\\
\toprule
\textbf{Test} & \textbf{Ce qu'il garantit} & \textbf{Résultat} \\
\midrule
\endfirsthead
\toprule
\textbf{Test} & \textbf{Ce qu'il garantit} & \textbf{Résultat} \\
\midrule
\endhead
\bottomrule
\endfoot
Le taux de complétion est calculé, non saisi
& Aucune valeur de complétion ne peut être écrite depuis l'API & OK \\
Une dimension non instruite ne compte pas comme un zéro
& Le score n'est pas mécaniquement dégradé par une absence de réponse & OK \\
Les moyennes d'axe ignorent les dimensions non instruites
& Les deux moyennes restent comparables entre campagnes & OK \\
Instruire une dimension fait progresser la complétion
& Le lien réponse $\rightarrow$ complétion est effectif & OK \\
Instruire une dimension modifie le score
& Le score est bien recalculé, non mis en cache & OK \\
Une dimension critique sous le seuil devient bloquante
& La règle bloquante / vigilance est appliquée & OK \\
Répondre au diagnostic exige la permission
& L'écriture est refusée sans \code{diagnostic.answer} & OK \\
La campagne d'un autre tenant est introuvable
& Réponse 404, pas 403 & OK \\
\end{longtable}

\subsection{Autorisation par rôle}

Quarante tests vérifient que chaque rôle fait ce qu'il doit faire \emph{et} échoue sur
ce qu'il ne doit pas faire. La seconde moitié est la plus importante : elle constitue la
preuve que les permissions sont réellement discriminantes.

\begin{longtable}{@{}L{7.4cm} L{5.2cm} C{1.5cm}@{}}
\caption{Extraits des tests d'autorisation}
\label{tab:tests-permissions}\\
\toprule
\textbf{Test} & \textbf{Ce qu'il garantit} & \textbf{Résultat} \\
\midrule
\endfirsthead
\toprule
\textbf{Test} & \textbf{Ce qu'il garantit} & \textbf{Résultat} \\
\midrule
\endhead
\bottomrule
\endfoot
Le catalogue ne contient aucun code inconnu
& Aucune permission orpheline dans la base & OK \\
Chaque modèle de rôle ne référence que des permissions existantes
& Pas de rôle pointant vers une permission supprimée & OK \\
Un administrateur d'entreprise ne peut pas s'octroyer le rôle de plateforme
& La frontière entre tenants n'est pas franchissable & OK \\
\ldots{} ni créer un rôle portant une permission de plateforme
& L'élévation de privilège par rôle personnalisé est bloquée & OK \\
Un rôle système n'est pas modifiable par un administrateur d'entreprise
& Les rôles livrés restent cohérents & OK \\
La Direction ne crée pas d'utilisateur
& Consulter et arbitrer n'est pas administrer & OK \\
Finance ne modifie pas la recommandation, ne pilote pas les projets
& Séparation des responsabilités & OK \\
L'AI Lead requalifie une opportunité, n'administre pas les comptes
& \code{recommendation.override} est exclusive à ce rôle & OK \\
Le Data Office ne gère pas le référentiel applicatif
& Data et architecture sont deux périmètres distincts & OK \\
Le PMO ne valide pas les bénéfices financiers
& La validation financière reste chez Finance et la Direction & OK \\
L'auditeur ne peut rien écrire, modifier ni supprimer
& Lecture seule effective sur tous les modules & OK \\
Le contributeur ne voit ni le portefeuille consolidé ni les données financières
& Le profil le plus restreint l'est réellement & OK \\
\end{longtable}

\subsection{Périmètre et étanchéité}

Dix tests couvrent le bornage par département et treize l'étanchéité entre
organisations. Ces derniers construisent une seconde organisation complète, avec ses
propres utilisateurs et rôles, et vérifient qu'aucune lecture, écriture, suppression,
affectation de département, attribution de rôle ni entrée d'audit ne traverse la
frontière. Le test final vérifie que toutes les listes du second locataire sont vides
des données du premier.

Un point mérite d'être explicité : un utilisateur d'une autre organisation qui appelle
une ressource existante reçoit \textbf{404}, non 403. Répondre 403 confirmerait
l'existence de la ressource, ce qui est déjà une fuite d'information. En revanche une
écriture hors périmètre \emph{à l'intérieur} de sa propre organisation reçoit bien 403,
puisque l'utilisateur voit l'objet mais n'a pas le droit d'agir dessus.

\subsection{Garde-fous structurels}

Quatre tests ne vérifient pas un comportement mais une propriété du code lui-même. Ils
parcourent la table de routage de l'application et contrôlent que :

\begin{itemize}
  \item aucun endpoint n'est exposé sans contrôle, à l'exception d'une liste explicite
    de routes publiques, chacune assortie d'un commentaire justifiant son ouverture ;
  \item aucune dépendance d'autorisation n'a été silencieusement transformée en
    paramètre de requête ou en corps de requête ;
  \item toutes les routes métier vivent sous le préfixe versionné.
\end{itemize}

Le deuxième de ces tests a été ajouté après un incident réel. Un import manquant avait
transformé une annotation de périmètre en paramètre de requête : FastAPI répondait 422
avec un message obscur, et surtout le contrôle de périmètre n'était plus appliqué. Le
test rend cette classe d'erreur impossible à réintroduire sans être détectée.

\section{Vérification des interfaces}

Les parcours suivants ont été exécutés dans un navigateur automatisé, avec capture à
chaque étape.

\begin{table}[H]
\centering
\caption{Parcours d'interface vérifiés}
\label{tab:tests-ui}
\small
\begin{tabular}{@{}L{5.6cm} L{6.4cm} C{1.8cm}@{}}
\toprule
\textbf{Parcours} & \textbf{Attendu} & \textbf{Résultat} \\
\midrule
Connexion par rôle
& Chaque rôle atterrit sur son écran d'accueil pertinent & OK \\
Navigation latérale
& Le menu ne contient que les écrans accessibles & OK \\
Panneau de diagnostic
& La saisie recalcule score, moyennes et complétion & OK \\
Recalcul de la priorisation
& Les curseurs modifient réellement le classement & OK \\
Création d'opportunité
& La qualification est déclenchée à l'enregistrement & OK \\
Export CSV du diagnostic
& Fichier téléchargé, accents corrects, résumé inclus & OK \\
Affichage mobile
& Aucun débordement horizontal jusqu'à 390~px & OK \\
\bottomrule
\end{tabular}
\end{table}

\section{Jeu de démonstration}

Les tests s'appuient sur des données créées à la volée ; la démonstration s'appuie sur
un jeu ensemencé, calibré pour être crédible et cohérent avec les maquettes.

\begin{table}[H]
\centering
\caption{Contenu du jeu de démonstration}
\label{tab:seed}
\small
\begin{tabular}{@{}L{6cm} C{2.4cm} L{5.4cm}@{}}
\toprule
\textbf{Élément} & \textbf{Volume} & \textbf{Remarque} \\
\midrule
Organisation & 1 & Laboratoire pharmaceutique \\
Sites & 4 & Tunis, Lyon, Casablanca, Paris \\
Départements & 8 & Du Commercial aux Affaires réglementaires \\
Utilisateurs & 13 & Au moins un par rôle fonctionnel \\
Réponses de diagnostic & 16 & Deux axes, avec preuves et confiances \\
Opportunités & 18 & Toutes qualifiées et scorées \\
Projets & 13 & Répartis sur les cinq phases \\
Entrées d'audit & 440 au seed & Volume qui croît ensuite à chaque écriture \\
\bottomrule
\end{tabular}
\end{table}

Les valeurs calculées à partir de ce jeu sont stables et reproductibles : readiness
global de 62,0 sur 100, moyennes d'axe de 2,81 et 3,12 sur~5, deux dimensions bloquantes
— gouvernance des données et infrastructure IA —, quatre premières priorités à 87, 81,
74 et 69~points, ROI de portefeuille de 183,9~\% et délai de retour moyen de 11~mois.

Ces valeurs ne sont écrites nulle part : elles résultent des formules appliquées aux
données ensemencées. C'est la raison pour laquelle le jeu a été calibré avec soin — une
démonstration dont les chiffres ne tiennent pas ensemble se voit immédiatement.

\begin{remarque}[Sur les comptes de démonstration]
Le mot de passe commun aux comptes du seed n'est pas écrit dans le code de l'interface.
Il est lu d'une variable d'environnement et n'est publié que par un endpoint d'aide qui
répond 404 dès que l'environnement n'est pas déclaré comme développement. Ces comptes
restent strictement destinés au développement et à la démonstration.
\end{remarque}

\section{Difficultés rencontrées}

Les difficultés ci-dessous ont été effectivement rencontrées pendant le développement et
sont documentées ici parce qu'elles ont conduit à des corrections vérifiables dans le
code.

\textbf{Une dimension ajoutée au référentiel dégradait tous les scores.} Les dimensions
non encore instruites étaient comptées comme des zéros. Ajouter une question au
questionnaire faisait donc chuter le readiness de toutes les campagnes existantes,
alors que rien n'avait changé sur le terrain. La notion de dimension instruite, adossée
à une confiance strictement positive, a résolu le problème : une dimension non mesurée
sort du calcul et pèse seulement sur le taux de complétion.

\textbf{Une gate franchie faisait reculer un projet.} Valider une gate déjà franchie
ramenait le projet à une phase antérieure. Un garde-fou répond désormais 409 avec un
message explicite selon que la gate est déjà franchie ou pas encore atteinte.

\textbf{Le recalcul des synergies écrasait les valeurs ensemencées.} La priorisation
recalculait automatiquement les synergies, ce qui modifiait les scores à chaque passage
et rendait la démonstration instable. Le recalcul est devenu explicite plutôt
qu'automatique.

\textbf{Le délai de retour ne concordait pas avec les hypothèses.} La formule initiale
rapportait le coût total de possession au gain mensuel, ce qui faisait porter au délai
de retour trois années de coûts récurrents non encore engagés. Elle a été corrigée pour
rapporter le seul investissement initial au gain annuel net d'exploitation.

\textbf{Un contrôle de périmètre avait disparu silencieusement.} Un import manquant
transformait une annotation en paramètre de requête ; FastAPI répondait 422 et le
contrôle n'était plus exécuté. Au-delà de la correction, un test structurel garantit
désormais qu'aucune dépendance ne peut subir ce sort sans être détectée.

\textbf{Un test passait seul et échouait en suite.} Un test d'interface laissait des
pondérations modifiées en base, ce qui faussait le test de priorisation exécuté ensuite.
Le montage de test rétablit désormais les pondérations par défaut avant et après.

\textbf{Le développement sous Windows a posé ses propres questions.} La documentation
d'installation décrivait l'activation d'environnement virtuel en syntaxe POSIX, ce qui
produisait une erreur sous PowerShell. Les deux fichiers de documentation donnent
désormais la commande adaptée à chaque terminal, avec une note sur la politique
d'exécution.
